\section{Nonregular languages and the pumping lemma}

\subsection{Languages that look nonregular}
We now know
\[
  \text{DFA}\;\equiv\;\text{NFA}\;\equiv\;\text{regular expressions}.
\]
Candidate nonregular languages:
\[
\begin{aligned}
  B&=\set{0^n1^n:n\geq 0},\\
  C&=\set{w\in\set{0,1}^*:w\text{ has equally many }0\text{s and }1\text{s}},\\
  F&=\set{ww:w\in\set{0,1}^*}.
\end{aligned}
\]
$B$ seems to require remembering how many $0$s have been seen. Finite memory
cannot do that for arbitrary $n$.

\begin{exercise}
Let $D$ be the language of strings over $\set{0,1}$ with equally many
occurrences of the substrings $01$ and $10$. Show that $D$ is regular.
(The language depends only on the first and last symbols.)
\end{exercise}

\begin{exercise}
Every finite language is regular. (Build a tree-like DFA, or a regular
expression that is a finite union of concrete words.)
\end{exercise}

\subsection{The pumping lemma}
If a DFA has $p$ states and a string of length at least $p$ is accepted, then
in the run some state repeats. The substring that takes the machine around
that loop can be repeated any number of times.

\begin{lemma}[Pumping lemma for regular languages]
If $A$ is regular, then there exists a pumping length $p\geq 1$ such that every
$s\in A$ with $|s|\geq p$ can be written $s=xyz$ satisfying
\begin{enumerate}
  \item $xy^iz\in A$ for every $i\geq 0$;
  \item $|y|>0$;
  \item $|xy|\leq p$.
\end{enumerate}
The strings $x$ and $z$ are allowed to be $\eps$.
\end{lemma}

\begin{proof}
Let $M=(Q,\Sigma,\transition,q_0,F)$ recognize $A$, and set $p=|Q|$. Let
$s=s_1s_2\cdots s_n\in A$ with $n\geq p$. Let $r_0,r_1,\ldots,r_n$ be the
states visited, so $r_0=q_0$, $r_{j}=\transition(r_{j-1},s_j)$, and
$r_n\in F$. The prefix sequence $r_0,\ldots,r_p$ has length $p+1$, so two of
these states coincide: take the first repeated state, at positions $i<j$
with $0\leq i<j\leq p$. Set
\[
  x=s_1\cdots s_i,\qquad
  y=s_{i+1}\cdots s_j,\qquad
  z=s_{j+1}\cdots s_n.
\]
Then $y$ takes $M$ from that state back to itself, so $xy^iz$ is accepted for
every $i\geq 0$. Also $|y|=j-i\geq 1$ and $|xy|=j\leq p$.
\end{proof}

\begin{center}
\begin{tikzpicture}[automaton,node distance=32mm]
  \node[state,initial] (q0) {$q_0$};
  \node[state,right=of q0] (qi) {$q$};
  \node[state,accepting,right=of qi] (qn) {$q_n$};
  \path[->]
    (q0) edge node[above]{$x$} (qi)
    (qi) edge[loop above] node{$y$} ()
    (qi) edge node[above]{$z$} (qn);
\end{tikzpicture}
\end{center}
The repeated state is $q=q_i=q_j$. Pumping $y$ stays in $q$ and still reaches
an accept state.

\subsection{How to prove a language is not regular}
Quantifier order:
\[
  \text{assume regular}
  \;\to\;
  \text{obtain }p
  \;\to\;
  \text{choose }s
  \;\to\;
  \text{defeat every }xyz
  \;\to\;
  \text{choose }i.
\]
You do not get to pick the decomposition; the adversary does. You must show
that \emph{every} $xyz$ obeying $|xy|\leq p$ and $|y|\geq 1$ fails for some
$i$.

\begin{example}[$B=\set{0^n1^n:n\geq 0}$]
Suppose $B$ is regular with pumping length $p$. Choose $s=0^p1^p$. Then
$|s|\geq p$. For any decomposition $s=xyz$ with $|xy|\leq p$ and $|y|\geq 1$,
the string $y$ lies entirely in the $0$s, so $y=0^k$ with $k\geq 1$. Pumping
with $i=0$ (or $i=2$) changes the number of $0$s but not the number of $1$s,
so $xy^iz\notin B$. Contradiction.
\end{example}

\begin{example}[$C=\set{w:\text{equal number of }0\text{s and }1\text{s}}$]
The same $s=0^p1^p$ works: $y$ is again a nonempty block of $0$s, and
$xy^2z$ has more $0$s than $1$s. (Choosing $s=(01)^p$ does \emph{not} work:
$y$ might be $01$, which pumps inside $C$.)
\end{example}

\begin{example}[$F=\set{ww:w\in\set{0,1}^*}$]
Take $s=0^p10^p1$. Then $|s|\geq p$. If $y$ lies in the first block of $0$s,
pumping breaks the two-copy structure. The Myhill--Nerode argument in the
next section is cleaner.
\end{example}

\begin{example}[$E=\set{0^i1^j:i\neq j}$]
A standard trick: if $E$ were regular then
$B=\overline{E}\cap 0^*1^*$ would be regular, but $B=\set{0^n1^n:n\geq 0}$.
Direct pumping on $0^{p}1^{p!+p}$ is another option (pumping $i=0$ is not
enough; the factorial gap prevents accidental equality).
\end{example}

\begin{exercise}
Show that $\set{1^{n^2}:n\geq 0}$ is not regular.
\end{exercise}
